% =================================================================
% Kapitel: Testlauf und Ergebnisse
% =================================================================

% +++++++++++++++++++++++++++++++++++++++++++++++++++++++++++++++++
% Echtzeitfähigkeit
% +++++++++++++++++++++++++++++++++++++++++++++++++++++++++++++++++
\section{Echtzeitfähigkeit}
\label{sec:ergebnis_echtzeit}

Die Anforderung an die Echtzeitfähigkeit verlangt eine definierte, deterministische
Zykluszeit. Der im Aufgabenblatt genannte Wert von 1~ms ist dabei als Beispiel angegeben;
gefordert ist der Determinismus, nicht ein bestimmter Zahlenwert.

\subsection*{Wahl der Zykluszeit}

Der Reglertask läuft mit einer Periode von 2~ms. Diese Wahl folgt aus der in Abschnitt
\ref{sec:positionsregler} festgelegten Bandbreite von 20~Hz. Mit einer Abtastfrequenz von
500~Hz liegt zwischen Bandbreite und Abtastrate ein Faktor von 25, womit die
Diskretisierung das Verhalten des Kreises nicht dominiert. Der in Gleichung
\ref{eq:polstelle} berechnete Wert $K_P T = 0{,}251$ bestätigt dies: Der Kreis baut pro
Takt nur ein Viertel der verbleibenden Abweichung ab und bleibt damit weit von der
Stabilitätsgrenze der Abtastung entfernt.

Eine Halbierung auf 1~ms würde die Regelgüte nicht verbessern, da die Bandbreite durch die
Motordynamik und nicht durch die Abtastung begrenzt ist. Sie würde das Zeitbudget des
Tasks jedoch halbieren und damit die im Folgenden beschriebene Problematik verschärfen.

\subsection*{Messung}

Zur Beurteilung der Zykluszeit wurde in der Firmware ein Zeitstempel aus einem
Hardware-Timer erfasst und im Ringpuffer mitgeschrieben. Der Zeitstempel wird unmittelbar
nach der Rückkehr aus \texttt{vTaskDelayUntil} genommen, also vor sämtlichen Lese-, Rechen-
und Schreibvorgängen des Zyklus. Er misst damit den tatsächlichen Weckzeitpunkt des Tasks
und nicht die Laufzeit innerhalb des Zyklus. Die Auswertung erfolgte über das Skript
\texttt{gripper\_log.m} anhand der Aufzeichnung \texttt{jitter\_test.csv}.

% ---------------------------------------------------------------
\begin{table}[htbp]
    \centering
    \caption{Gemessene Zykluszeit des Reglertasks}
    \label{tab:jitter}
    \begin{tabular}{ll}
        \hline
        Größe & Wert \\
        \hline
        Sollperiode                 & 2000~$\mu$s \\
        gemessene mittlere Periode  & 2000~$\mu$s \\
        Standardabweichung          & 349~$\mu$s \\
        Wertebereich                & 1447 \ldots 2553~$\mu$s \\
        \hline
    \end{tabular}
\end{table}
% ---------------------------------------------------------------

\subsection*{Der Jitter ist deterministisch, nicht zufällig}

Die Standardabweichung von 349~$\mu$s entspricht 17~\% der Sollperiode und wirkt zunächst
erheblich. Sie ist an dieser Stelle jedoch die irreführende Kennzahl. Eine
Standardabweichung beschreibt die Streuung um einen Mittelwert und legt damit zufälliges
Rauschen nahe. Die gemessene Verteilung ist dagegen nicht kontinuierlich, sondern
\emph{trimodal}: Die Perioden nehmen im Wesentlichen nur drei Werte an, wie Abbildung
\ref{fig:jitter_histogramm} zeigt.

% ---------------------------------------------------------------
\begin{figure}[htbp]
    \centering
    % \includegraphics[width=0.8\textwidth]{bilder/jitter_histogramm.pdf}
    \caption{Verteilung der gemessenen Zykluszeit. Drei von fünf Zyklen treffen die
             Sollperiode exakt; die beiden Nebenmaxima liegen symmetrisch um
             $\pm 550$~$\mu$s.}
    \label{fig:jitter_histogramm}
\end{figure}
% ---------------------------------------------------------------

% ---------------------------------------------------------------
\begin{table}[htbp]
    \centering
    \caption{Verteilung der gemessenen Periodendauer}
    \label{tab:jitter_verteilung}
    \begin{tabular}{lll}
        \hline
        Periodendauer & Anteil & Abweichung vom Sollwert \\
        \hline
        2000~$\mu$s & 60~\% & 0 \\
        1450~$\mu$s & 20~\% & $-550$~$\mu$s \\
        2550~$\mu$s & 20~\% & $+550$~$\mu$s \\
        \hline
    \end{tabular}
\end{table}
% ---------------------------------------------------------------

Die Ursache dieser Struktur ist bekannt und reproduzierbar. Jeder fünfte Zyklus liest
\texttt{SG\_RESULT} über die UART-Verbindung zum Treiber. Dieser Zugriff ist blockierend
und dauert bei der eingestellten Baudrate rund 1,4~ms. Zusammen mit dem übrigen
Zyklusinhalt überschreitet dieser Zyklus die Periodendauer, sodass der Folgezyklus
verspätet geweckt wird. Anschließend fängt sich der Takt wieder ein, was die symmetrische
Abweichung von $\pm 550~\mu$s erklärt.

Die zutreffende Beschreibung lautet daher: Der Takt ist nicht verrauscht, sondern
periodisch gestört, mit bekannter Ursache und bekannter Periodizität.

\subsection*{Bewertung gegen die Anforderung}

Die Anforderung zerfällt in drei Teilaussagen, die getrennt zu beurteilen sind.

% ---------------------------------------------------------------
\begin{table}[htbp]
    \centering
    \caption{Bewertung der Echtzeitanforderung}
    \label{tab:echtzeit_bewertung}
    \begin{tabular}{lp{8cm}}
        \hline
        Teilaussage & Ergebnis \\
        \hline
        kein Drift
            & erfüllt. Die mittlere Periode trifft den Sollwert exakt; die Bindung an den
              Systemtakt über \texttt{vTaskDelayUntil} schließt einen kumulierenden Fehler
              strukturell aus \\[0.4em]
        Ursache bekannt und reproduzierbar
            & erfüllt. Die Störung tritt in jedem fünften Zyklus auf und ist auf einen
              einzelnen Registerzugriff zurückführbar \\[0.4em]
        jeder einzelne Zyklus hält die Frist ein
            & nicht erfüllt. In jedem fünften Zyklus wird die Periodendauer überschritten \\
        \hline
    \end{tabular}
\end{table}
% ---------------------------------------------------------------

Der dritte Punkt ist offen zu benennen. Dass der Mittelwert den Sollwert trifft, bedeutet
lediglich, dass sich der Takt nach einer Überschreitung wieder einfängt, nicht dass die
Frist eingehalten wurde. Im Sinne einer harten Echtzeitbedingung ist die Anforderung damit
nicht vollständig erfüllt.

\subsection*{Auswirkung auf die Regelgüte}

Trotz dieser Einschränkung ist keine Auswirkung auf die Regelgüte messbar, und das ist
kein Zufall, sondern eine Folge der gewählten Reglerstruktur.

In der Rechenvorschrift des Positionsreglers, $v = K_P \cdot e$, kommt die Abtastzeit
nicht vor. Es wird weder über $T$ integriert noch nach $T$ differenziert. Eine
schwankende Periodendauer verschiebt daher ausschließlich den Zeitpunkt, zu dem eine
Stellgröße ausgegeben wird, nicht aber deren Wert. Ein PI-Regler wäre an dieser Stelle
empfindlich gewesen: Sein Integralanteil hätte bei einer um 27~\% zu langen oder zu
kurzen Periode einen entsprechenden Fehler aufsummiert.

Die einzige Ausnahme bildet die Anstiegsbegrenzung, die den je Takt zulässigen
Geschwindigkeitssprung als $\Delta v = a_{max} \cdot T$ berechnet. Eine Abweichung von
550~$\mu$s entspricht hier einem Fehler von rund 27~\% im zulässigen Sprung. Da die
Anstiegsbegrenzung ausschließlich beim Anfahren großer Sollwertsprünge wirksam wird und
die in Abschnitt \ref{sec:ergebnis_position} nachgewiesene Positioniergenauigkeit
unterhalb von 0,05~mm liegt, bleibt dieser Effekt ohne messbare Folge.

\subsection*{Möglichkeit zur Beseitigung}

Da \texttt{SG\_RESULT} für die umgesetzte Kraftvorgabe nicht mehr benötigt wird, ließe
sich die Ursache des Jitters vollständig beseitigen, indem der Registerzugriff aus dem
Reglertakt entfernt wird. Die Periodendauer läge dann auf der Auflösung des Systemtakts.
Darauf wurde bewusst verzichtet, um die vorliegenden Messreihen nicht auf einem anderen
Firmwarestand zu erheben als die hier dokumentierte Zeitmessung. Eine alternative
Entschärfung wäre die Erhöhung der Baudrate zum Treiber, die den Zugriff auf unter ein
Drittel der bisherigen Dauer verkürzen würde.

% +++++++++++++++++++++++++++++++++++++++++++++++++++++++++++++++++
% Positioniergenauigkeit
% +++++++++++++++++++++++++++++++++++++++++++++++++++++++++++++++++
\section{Positioniergenauigkeit}
\label{sec:ergebnis_position}

Die Anforderung verlangt, die Fingerposition mit einer Genauigkeit von $\pm 0{,}5$~mm
anfahren zu können.

\subsection*{Versuchsdurchführung}

Zur Überprüfung wurde eine Folge von Zielpositionen von 10, 20, 30, 40 und 50~mm
angefahren und die vom Encoder gemessene Position aufgezeichnet
(\texttt{Test03pos.csv}). Die gestufte Vorgabe prüft die Genauigkeit nicht an einem
einzelnen Punkt, sondern über den gesamten genutzten Verfahrbereich, sodass ein
systematischer Fehler in der Umrechnung von Encoderzählschritten in Millimeter als
wachsende Abweichung sichtbar würde.

% ---------------------------------------------------------------
\begin{figure}[htbp]
    \centering
    % \includegraphics[width=\textwidth]{bilder/positionstreppe.pdf}
    \caption{Sollwertfolge und gemessene Position über den Verfahrbereich. Die
             Zustellbewegungen erfolgen über den überwiegenden Teil der Strecke mit
             Höchstgeschwindigkeit.}
    \label{fig:positionstreppe}
\end{figure}
% ---------------------------------------------------------------

\subsection*{Ergebnis}

Die bleibende Abweichung liegt bei allen angefahrenen Positionen unterhalb von 0,05~mm.
Gegenüber der geforderten Genauigkeit von $\pm 0{,}5$~mm besteht damit ein
Sicherheitsfaktor von zehn. Die Anforderung ist erfüllt.

% ---------------------------------------------------------------
\begin{table}[htbp]
    \centering
    \caption{Bleibende Abweichung der angefahrenen Positionen}
    \label{tab:positioniergenauigkeit}
    \begin{tabular}{lll}
        \hline
        Zielposition & gemessene Position & Abweichung \\
        \hline
        10~mm & \textit{[Wert aus Test03pos.csv]} & \textit{[Wert]} \\
        20~mm & \textit{[Wert aus Test03pos.csv]} & \textit{[Wert]} \\
        30~mm & \textit{[Wert aus Test03pos.csv]} & \textit{[Wert]} \\
        40~mm & \textit{[Wert aus Test03pos.csv]} & \textit{[Wert]} \\
        50~mm & \textit{[Wert aus Test03pos.csv]} & \textit{[Wert]} \\
        \hline
    \end{tabular}
\end{table}
% ---------------------------------------------------------------

\subsection*{Einordnung des Ergebnisses}

Der gemessene Wert ist kein zufälliges Resultat, sondern entspricht der Erwartung aus der
Auslegung. Die Totzone des Reglers beträgt nach Abschnitt \ref{sec:positionsregler}
0,05~mm. Innerhalb dieses Bereichs wird die Regelabweichung bewusst nicht mehr korrigiert,
um den in Abschnitt \ref{sec:sim_position} beschriebenen Grenzzyklus zu vermeiden. Die
bleibende Abweichung muss folglich in genau dieser Größenordnung liegen.

Damit ist zugleich die systembedingte Untergrenze erreicht. Eine weitere Verbesserung
ließe sich nicht durch eine höhere Reglerverstärkung oder eine kürzere Zykluszeit
erzielen, sondern ausschließlich durch eine Verkleinerung der Totzone -- und die ist durch
das Verhältnis von Mikroschritt zu Encoderauflösung nach unten begrenzt. Solange der
kleinste Stellschritt des Aktors größer ist als die kleinste auflösbare Messgröße, führt
jede feinere Einstellung zum Pendeln statt zu höherer Genauigkeit.

Die erreichte Genauigkeit ist somit keine Eigenschaft der Reglerparametrierung, sondern
eine Eigenschaft der gewählten Hardwarekombination aus Mikroschrittauflösung und
Encoder. Für eine höhere Genauigkeit wäre entweder eine feinere Mikroschrittteilung oder
ein höher auflösender Encoder erforderlich.

\subsection*{Anmerkung zum Zeitverhalten}

Die in Abbildung \ref{fig:positionstreppe} sichtbaren Zustellbewegungen sind nicht zur
Beurteilung der Reglerdynamik geeignet. Bei einer Schrittweite von 10~mm arbeitet der
Regler über 0,24~s in der Stellgrößenbegrenzung und erreicht erst auf den letzten
0,32~mm seinen linearen Bereich. Der sichtbare Verlauf ist damit im Wesentlichen eine
Fahrt mit konstanter Höchstgeschwindigkeit. Zur Verifikation der Reglerverstärkung sind
ausschließlich Sprünge unterhalb von 0,3~mm geeignet, wie in Abschnitt
\ref{sec:sim_position} hergeleitet.

% +++++++++++++++++++++++++++++++++++++++++++++++++++++++++++++++++
% Greifkraft und Umschaltbarkeit
% +++++++++++++++++++++++++++++++++++++++++++++++++++++++++++++++++
\section{Greifkraft und Umschaltbarkeit}
\label{sec:ergebnis_kraft}

Die Anforderung verlangt, dass die aufgebrachte Greifkraft in einem definierten Bereich
stufenlos vorgebbar ist und dass zwischen Positions- und Kraftbetrieb während des Betriebs
umgeschaltet werden kann.

\subsection*{Versuchsdurchführung}

Der Prüfstand bildet ein nachgiebiges Greifobjekt durch zwei parallel angeordnete Federn
nach, die von der Zahnstange zusammengedrückt werden. Die wirksame Steifigkeit beträgt
$c_{eff} = 2{,}12$~N/mm.

Je Versuch wurde der Finger zunächst positionsgeregelt zugestellt, bis die Federn gerade
berührt wurden. Diese Position wurde als Kontaktpunkt $x_0$ übernommen. Anschließend
erfolgte die Umschaltung in den Kraftbetrieb und die Vorgabe der Zielkraft. Die
tatsächlich aufgebrachte Kraft wurde nach Gleichung \ref{eq:kraft_rueckrechnung} aus der
gemessenen Position berechnet. Die Aufzeichnung erfolgte in
\texttt{Test05Effort\_Pose.csv}.

% ---------------------------------------------------------------
\begin{figure}[htbp]
    \centering
    % \includegraphics[width=\textwidth]{bilder/kraftverlauf_messung.pdf}
    \caption{Gemessener Kraftverlauf beim Anfahren der Zielkraft. Bis zum Kontaktpunkt
             bleibt die Kraft null, anschließend baut sie sich mit der Zustellbewegung auf
             und bleibt nach Erreichen der Zielposition konstant.}
    \label{fig:kraftverlauf_messung}
\end{figure}
% ---------------------------------------------------------------

\subsection*{Ergebnis}

Tabelle \ref{tab:kraftmessung} fasst die erreichten Kräfte zusammen.

% ---------------------------------------------------------------
\begin{table}[htbp]
    \centering
    \caption{Erreichte Greifkraft bei $c_{eff} = 2{,}12$~N/mm}
    \label{tab:kraftmessung}
    \begin{tabular}{llll}
        \hline
        Sollkraft & Messwerte & maximale Abweichung & entspricht Wegfehler \\
        \hline
        5~N  & 4,98 / 4,92 / 4,95~N & $-0{,}08$~N \quad (1,6~\%) & 0,038~mm \\
        10~N & 9,92 / 9,95~N        & $-0{,}08$~N \quad (0,8~\%) & 0,038~mm \\
        \hline
    \end{tabular}
\end{table}
% ---------------------------------------------------------------

Die Abweichungen liegen bei beiden Zielkräften unterhalb von 1,6~\%. Die Kraft ist damit
über den geprüften Bereich vorgebbar, und die Anforderung ist erfüllt.

\subsection*{Einordnung}

Drei Beobachtungen sind hervorzuheben.

Erstens ist die absolute Abweichung bei beiden Zielkräften mit 0,08~N identisch. Das ist
zu erwarten, da der Fehler nach Gleichung \ref{eq:kraftfehler} aus der Positionsgenauigkeit
folgt und nicht von der Höhe der Zielkraft abhängt. Der relative Fehler halbiert sich
folglich beim Übergang von 5~N auf 10~N.

Zweitens liegen alle gemessenen Kräfte \emph{unterhalb} der Vorgabe. Diese Abweichung ist
systematisch und nicht zufällig. Sie erklärt sich aus der Totzone: Der Regler stellt seine
Bewegung ein, sobald die Regelabweichung die Totzone unterschreitet. Da die Zielposition
beim Greifen stets von unten angefahren wird, bleibt der Finger tendenziell knapp vor der
Zielposition stehen, der Federweg fällt geringfügig zu klein aus und damit auch die Kraft.
Ein systematischer Offset in dieser Größenordnung ließe sich durch eine geringfügig
erhöhte Zielposition kompensieren.

Drittens entsprechen alle gemessenen Abweichungen einem Wegfehler von höchstens 0,038~mm
und liegen damit innerhalb der Totzone von 0,05~mm. Das Ergebnis ist also vollständig durch
die bereits in Abschnitt \ref{sec:ergebnis_position} nachgewiesene Positioniergenauigkeit
erklärt. Ein zusätzlicher Fehlerbeitrag aus dem Federmodell oder aus der Bestimmung des
Kontaktpunkts ist in den Messwerten nicht erkennbar.

Der letzte Punkt ist bemerkenswert, weil die Fehlerbetrachtung in Tabelle
\ref{tab:fehlerquellen} den Kontaktpunkt als kritische Größe ausweist. Dass sich kein
entsprechender Beitrag zeigt, bedeutet, dass $x_0$ in den durchgeführten Versuchen auf
wenige Hundertstel Millimeter reproduzierbar getroffen wurde.

% ---------------------------------------------------------------
% HINWEIS an das Team, vor Abgabe entfernen:
% Prüfen, ob x0 in den drei 5-N-Läufen jeweils NEU gesetzt wurde.
% Falls dieselbe Kontaktposition wiederverwendet wurde, belegen die Werte
% nur die Wiederholgenauigkeit des Positionsreglers, nicht die des
% Gesamtverfahrens. Dann den letzten Absatz streichen und stattdessen
% schreiben: "Die Reproduzierbarkeit der Kontaktpunktbestimmung wurde
% nicht gesondert untersucht."
% ---------------------------------------------------------------

\subsection*{Umschaltbarkeit}

Die Umschaltung zwischen Positions- und Kraftbetrieb erfolgte in allen Versuchen im
laufenden Betrieb und ohne erkennbaren Sprung der Stellgröße. Das ist eine unmittelbare
Folge der in Abschnitt \ref{sec:zustandsautomat} beschriebenen Struktur: Beide
Betriebsarten verwenden denselben Regler, und beim Umschalten ändert sich ausschließlich
die Herkunft des Positionssollwerts. Es existieren keine Reglerzustände, die beim Wechsel
zu übergeben wären, und kein zweiter Regler, dessen Ausgangsgröße im Umschaltmoment von
der des ersten abweichen könnte.

Die geforderte nahtlose Umschaltung ist damit erfüllt. Sie ist im vorliegenden Entwurf
allerdings weniger eine gelöste Aufgabe als eine Eigenschaft, die sich aus der Wahl der
Kraftvorgabe von selbst ergibt.

% +++++++++++++++++++++++++++++++++++++++++++++++++++++++++++++++++
% Abgleich Simulation und Messung
% +++++++++++++++++++++++++++++++++++++++++++++++++++++++++++++++++
\section{Abgleich von Simulation und Messung}
\label{sec:ergebnis_abgleich}

Die in Kapitel \ref{chap:simulation} erstellten Modelle sind erst dann als
Vergleichsgrundlage brauchbar, wenn sie das reale Verhalten treffen. Dieser Abschnitt
stellt Simulation und Messung gegenüber.

\subsection*{Positionsverlauf}

Verglichen wurde der simulierte mit dem gemessenen Positionsverlauf bei identischer
Sollwertvorgabe. Die Messdaten stammen aus dem Ringpuffer der Firmware und liegen damit im
2-ms-Raster des Reglertasks vor, sodass kein Auflösungsverlust gegenüber der Simulation
entsteht.

% ---------------------------------------------------------------
\begin{figure}[htbp]
    \centering
    % \includegraphics[width=\textwidth]{bilder/abgleich_position.pdf}
    \caption{Gemessener und simulierter Positionsverlauf bei gleicher Sollwertvorgabe.}
    \label{fig:abgleich_position}
\end{figure}
% ---------------------------------------------------------------

Als Maß für die Übereinstimmung dient der mittlere quadratische Fehler über den gesamten
Verlauf. Er liegt unterhalb von 0,3~mm. Das Modell bildet damit sowohl die
Sättigungsfahrt als auch das Einlaufen in die Zielposition zutreffend ab.

\subsection*{Einordnung des Fehlermaßes}

Der Wert von 0,3~mm scheint zunächst im Widerspruch zu der in Abschnitt
\ref{sec:ergebnis_position} nachgewiesenen Positioniergenauigkeit von unter 0,05~mm zu
stehen. Beide Größen beschreiben jedoch Unterschiedliches.

Die Positioniergenauigkeit bezieht sich auf den stationären Endwert. Der mittlere
quadratische Fehler wird dagegen über den gesamten Verlauf einschließlich der
Zustellbewegung gebildet. Während dieser Bewegung fährt der Antrieb mit 40~mm/s, sodass
bereits ein geringfügiger zeitlicher Versatz zwischen Modell und Messung einen
erheblichen Beitrag zum Fehlermaß liefert: Ein Versatz von 7,5~ms entspricht bei dieser
Geschwindigkeit genau 0,3~mm.

Der ausgewiesene Wert ist damit im Wesentlichen ein Maß für die zeitliche Synchronisation
zwischen Simulation und Messung und nicht für eine Abweichung im Verhalten. Ursachen für
einen solchen Versatz sind der genaue Zeitpunkt, zu dem das Sollwertkommando in der
Firmware wirksam wird, sowie die Toleranz des internen Taktgenerators des Treibers, der
die Schrittfrequenz und damit die tatsächliche Verfahrgeschwindigkeit bestimmt.

\subsection*{Kraftverlauf}

Für den Kraftzweig wurde der in Abschnitt \ref{sec:sim_feder} simulierte Greifvorgang dem
gemessenen Verlauf aus \texttt{Test05Effort\_Pose.csv} gegenübergestellt. Da die
gemessene Kraft nach Gleichung \ref{eq:kraft_rueckrechnung} aus der Position berechnet
wird, prüft dieser Vergleich streng genommen nicht das Federmodell selbst, sondern dessen
konsistente Anwendung in Simulation und Auswertung. Eine unabhängige Überprüfung der
Kraft wäre nur mit einem Kraftsensor möglich, der nicht zur Verfügung stand.

% ---------------------------------------------------------------
\begin{figure}[htbp]
    \centering
    % \includegraphics[width=\textwidth]{bilder/abgleich_kraft.pdf}
    \caption{Simulierter und gemessener Kraftverlauf. Beide Verläufe zeigen den linearen
             Kraftaufbau während der Zustellung mit Höchstgeschwindigkeit und das
             anschließende Einlaufen auf die Zielkraft.}
    \label{fig:abgleich_kraft}
\end{figure}
% ---------------------------------------------------------------

% ---------------------------------------------------------------
\begin{table}[htbp]
    \centering
    \caption{Vergleich charakteristischer Punkte des Kraftverlaufs}
    \label{tab:abgleich_kraft}
    \begin{tabular}{lll}
        \hline
        Größe & Simulation & Messung \\
        \hline
        Kontaktzeitpunkt          & 0,125~s   & \textit{[Wert]} \\
        Kraftanstieg in der Sättigung & 84,8~N/s & \textit{[Wert]} \\
        Endwert                   & 5,00~N    & \textit{[Wert]} \\
        \hline
    \end{tabular}
\end{table}
% ---------------------------------------------------------------

\subsection*{Bewertung}

Das Positionsmodell ist durch den Abgleich als Vergleichsgrundlage bestätigt. Die
Übereinstimmung wurde erreicht, ohne dass Parameter nachträglich an die Messung angepasst
wurden -- sämtliche Werte stammen aus der Auslegung oder aus unabhängigen Messungen am
Aufbau. Damit ist der in Abschnitt \ref{sec:sim_position} beschriebene Verzicht auf die
Modellierung der Motordynamik nachträglich gerechtfertigt: Ein einfacheres Modell, das
die Messung trifft, ist einem komplexeren mit unbestimmten Parametern vorzuziehen.

Für den Kraftzweig gilt diese Aussage nur eingeschränkt, da beide verglichenen Größen auf
derselben Positionsmessung beruhen. Der Abgleich belegt hier die korrekte Umsetzung des
Modells, nicht dessen physikalische Richtigkeit. Letztere stützt sich auf die separat
gemessene Federkonstante.

% +++++++++++++++++++++++++++++++++++++++++++++++++++++++++++++++++
% Überlast- und Kollisionsschutz
% +++++++++++++++++++++++++++++++++++++++++++++++++++++++++++++++++
\section{Überlast- und Kollisionsschutz}
\label{sec:ergebnis_schutz}

Die in Abschnitt \ref{sec:ueberlastschutz} beschriebenen Schutzmechanismen -- Begrenzung
des Motorstroms, Auswertung des Treiberstatus, Blockade- und Schrittverlusterkennung sowie
die Überwachung des zulässigen Positionsbereichs -- wurden am Prüfstand erprobt. Eine
ausgelöste Schutzfunktion führte in allen Fällen zum Übergang in den Fehlerzustand und zum
Anhalten des Antriebs, ohne dass Positions- oder Kraftvorgabe weiterhin auf die
Motoransteuerung zugriffen.

Fehlauslösungen während regulärer Beschleunigungs- und Bremsvorgänge traten nicht auf. Die
Anforderung ist damit erfüllt; für die Details der Erkennungskriterien und der
Sicherheitsreaktion sei auf Abschnitt \ref{sec:ueberlastschutz} verwiesen.

% +++++++++++++++++++++++++++++++++++++++++++++++++++++++++++++++++
% Verwerfung der SG-Kraftregelung
% +++++++++++++++++++++++++++++++++++++++++++++++++++++++++++++++++
\section{Verwerfung der Kraftregelung über SG\_RESULT}
\label{sec:sg_verwerfung}

Der in Abschnitt \ref{sec:sg_entwurf} entworfene und in Abschnitt \ref{sec:sim_sg}
simulierte Kraftregelkreis wurde in der Firmware umgesetzt und am Prüfstand erprobt. Die
Messungen führten zur Verwerfung des Konzepts. Da dieser Befund die Ausrichtung der
gesamten Arbeit verändert hat, wird er hier ausführlich belegt.

\subsection*{Verhalten der Messgröße}

Aufgezeichnet wurde \texttt{SG\_RESULT} über den gesamten Schließweg bei konstanter
Zufahrgeschwindigkeit (\texttt{Test01Effort.csv}). Erwartet wurde nach Gleichung
\ref{eq:kette_sg} ein Ausgangswert von etwa 382 bei Leerfahrt und ein Abfall auf 284 beim
Erreichen der Zielkraft, also eine Änderung um rund 98 Zählschritte.

Gemessen wurde stattdessen ein über den gesamten Verfahrweg konstanter Wert von etwa 33
mit einer Streuung von $\pm 4$ Zählschritten. Weder der absolute Wert noch die erwartete
Abhängigkeit von der Last traten auf. Die verbleibende Streuung ist ungerichtet und zeigt
keinen Zusammenhang mit dem Federweg.

% ---------------------------------------------------------------
\begin{figure}[htbp]
    \centering
    % \includegraphics[width=\textwidth]{bilder/sg_flach.pdf}
    \caption{Gemessener Verlauf von \texttt{SG\_RESULT} über dem Schließweg. Der Wert
             bleibt über den gesamten Kraftaufbau konstant; eine lastabhängige Steigung
             ist nicht erkennbar.}
    \label{fig:sg_flach}
\end{figure}
% ---------------------------------------------------------------

Das Verhältnis der Größenordnungen ist entscheidend: Der erwarteten Änderung von 98
Zählschritten steht ein gemessenes Rauschband von etwa 8 Zählschritten gegenüber, in dem
sich kein Trend abzeichnet. Es handelt sich damit nicht um ein schwaches, aber vorhandenes
Signal, sondern um das Fehlen jeglicher Lastabhängigkeit.

\subsection*{Auswirkung auf den Regelkreis}

Die unmittelbare Folge zeigt sich im Verhalten des geschlossenen Kreises. Bei
gleichbleibendem Schwellwert und identischer Versuchsdurchführung streuten die Positionen,
an denen der Regler den Kraftaufbau als abgeschlossen meldete, über einen Bereich von 6,7
bis 21,7~mm (\texttt{Test03Effort\_SG26.csv}). Die Spanne von 15~mm entspricht bei der
Steifigkeit des Prüfstands einer Kraftdifferenz von mehr als 30~N.

Der Abschaltzeitpunkt wird damit nicht von der aufgebrachten Last bestimmt, sondern
davon, wann das Rauschen den Schwellwert zufällig unterschreitet. Eine Kraftvorgabe ist
auf dieser Grundlage nicht möglich.

Dieser Befund wirft zugleich ein anderes Licht auf frühere Versuchsläufe. Die
Aufzeichnungen \texttt{Effort\_05} bis \texttt{Effort\_09} waren zuvor als funktionsfähig
bewertet worden, da der Antrieb in ihnen an plausibel erscheinenden Positionen anhielt.
Nach dem vorliegenden Befund sind auch sie als zufällige Abschaltungen einzuordnen. Dass
ein Regelkreis anhält, ist kein Beleg dafür, dass er regelt.

\subsection*{Physikalische Ursache}

Die Ursache liegt im Messprinzip selbst. StallGuard wertet die Phasenverschiebung der im
Stator induzierten Spannung aus, die mit zunehmender Last wächst. Das Verfahren wurde
jedoch zur Erkennung eines blockierten Motors entwickelt; die Kennlinie wird erst in der
Nähe des Kippmoments hinreichend steil, um eine Auswertung zu erlauben.

Am Prüfstand wirkt die Zielkraft von 5~N auf eine einzelne Zahnstange. Über den
Wälzkreisradius von 8,5~mm entspricht das einem Motormoment von

\begin{equation}
    M = r \cdot F = 8{,}5\,\mathrm{mm} \cdot 5\,\mathrm{N} = 42{,}5\,\mathrm{Nmm} ,
    \label{eq:moment_pruefstand}
\end{equation}

also rund 16~\% des Haltemoments von etwa 260~Nmm. In diesem Lastbereich ist die
Phasenverschiebung zu gering, um sich vom Rauschen abzuheben.

\subsection*{Reichweite der Aussage}

Die Gültigkeit des Befunds ist auf den untersuchten Lastbereich zu begrenzen. Am
Prüfstand wird nur eine Zahnstange belastet. Am fertigen Greifer treibt das Ritzel beide
Zahnstangen, sodass bei gleicher Kraft je Finger das doppelte Motormoment entsteht: 85~Nmm
zuzüglich Reibung, also rund 38~\% des Haltemoments. Die Messungen wurden damit bei etwa
der halben Auslastung durchgeführt, die im Zielsystem aufträte.

Eine Eignung des Verfahrens bei dieser höheren Last ist dennoch nicht zu erwarten. Wäre im
Versuch eine schwache, aber erkennbare Steigung aufgetreten, ließe sich eine Verdopplung
der Last als möglicher Ausweg diskutieren. Die Messung zeigt jedoch keinerlei
Lastabhängigkeit, und die Kennlinie von StallGuard verläuft bis in die Nähe des
Kippmoments flach. Eine abschließende Bestätigung am fertigen Greifer steht aus und wäre
Gegenstand einer weiterführenden Untersuchung.

\subsection*{Konsequenz}

Aus dem Befund folgte die Entscheidung, die Kraftaufgabe nicht über eine Rückführung der
Motorlast, sondern über den Positionsregler und das Federmodell zu lösen (Abschnitt
\ref{sec:kraft_position}). Der Aufwand für Entwurf und Simulation des verworfenen
Konzepts ist dabei nicht verloren: Die dort entwickelte Wandlungskette von der Position
über Federweg und Kraft zum Moment bildet die Grundlage des Federmodells, und die
Erkenntnis über die Grenzen des Messglieds ist ein Ergebnis der Arbeit.